\section{How prompt caching works, and what a probe showed}\label{sec:caching}

\paragraph{The idea.} Every request an agent sends repeats the whole conversation so far: the tool
definitions, the system prompt, and every earlier turn. Re-processing all of it each time would be slow
and expensive, so the provider keeps a \term{prompt cache}. The first time a prefix is sent it is
\term{written} to the cache, at a premium. Later requests that start with the same prefix \term{read} it
back at a fraction of the normal price. \Cref{tab:prices} gives the prices used throughout this report.

\begin{table}[htbp]
\centering\small
\caption{Price per million tokens (US dollars), by model and token class. The same table priced every
request in the campaign and in the probe.}
\label{tab:prices}
\begin{tabular*}{0.8\textwidth}{@{\extracolsep{\fill}}l r r r r@{}}
\toprule
Model & uncached input & cache read & cache write & output \\
\midrule
Fable 5.1 (host) & \PriceFableIn & \PriceFableRead & \PriceFableWrite & \PriceFableOut \\
Opus 5.5 (host) & \PriceOpusIn & \PriceOpusRead & \PriceOpusWrite & \PriceOpusOut \\
Sonnet 5 (cheaper model) & \PriceSonnetIn & \PriceSonnetRead & \PriceSonnetWrite & \PriceSonnetOut \\
\bottomrule
\end{tabular*}
\end{table}

Two things in this table drive the whole report. Fable is expensive on every class, so moving work to
Sonnet can save a lot. Opus is cheap to \emph{read} from cache: Sonnet reads cost \SonnetOpusReadPriceX\X\
Opus reads, while Sonnet writes cost only \SonnetOpusWritePriceX\X\ Opus writes. In a long conversation
that is mostly cache reads, Sonnet can cost more than Opus for the same tokens.

\paragraph{Where the prices come from, and how much the Opus result depends on them.}\label{sec:breakeven}
The table is the bundle's own price table (\code{savings.DEFAULT\_RATES}), last verified against the
provider's public price list on \RatesVerified. The provider's per-request cost field is computed from the
same table, so the exact match between our recomputed costs and the provider's (\cref{sec:design}) checks
that tokens were counted correctly, not that the prices are right. The Opus result rests mainly on one
number, Opus's cache-read price. \Cref{tab:breakeven} re-prices the recorded tokens of the test split,
changing only that rate and keeping everything else, including every token count, as measured.

\begin{table}[htbp]
\centering\small
\caption{Opus-host cost ratios (test split, tools-normalized) if Opus cache reads cost more than the table's
\$\PriceOpusRead\ per million tokens. Only that rate changes; every token count is as recorded.}
\label{tab:breakeven}
\begin{tabular*}{\textwidth}{@{\extracolsep{\fill}}l r r r r r r@{}}
\toprule
 & \multicolumn{3}{c}{Pair reading (co-primary)} & \multicolumn{3}{c}{Arm reading (co-primary, unfiltered)} \\
\cmidrule(lr){2-4}\cmidrule(l){5-7}
Opus cache-read price (per M) & sticky & shipped & sonnet & sticky & shipped & sonnet \\
\midrule
\$0.20 (table) & 1.25 & 1.38 & 1.43 & 1.20 & 1.32 & 1.46 \\
\$0.25 & 1.17 & 1.29 & 1.33 & 1.12 & 1.24 & 1.36 \\
\$0.30 & 1.09 & 1.22 & 1.25 & 1.05 & 1.16 & 1.27 \\
\$0.35 & 1.03 & 1.15 & 1.17 & 0.99 & 1.10 & 1.19 \\
\$0.40 & 0.98 & 1.09 & 1.11 & 0.94 & 1.04 & 1.13 \\
\bottomrule
\end{tabular*}

\end{table}
\howtoreadtable{The first row reproduces the confirmatory ratios. Reading down a column shows how the extra
cost of routing shrinks as Opus reads get dearer. Sticky reaches parity (ratio 1) at about
\$\BEStickyPair\ per million tokens in the pair reading and \$\BEStickyArm\ in the arm reading, roughly
\BEReadMultiple\X\ the table price; shipped at \$\BEShippedPair\ and plain Sonnet at \$\BESonnetPair\ (pair
reading). On the raw cost basis with the arm reading the same exercise gives sticky \BERawStickyThirty,
\BERawStickyThirtyFive\ and \BERawStickyForty\ at \$\RateThirty, \$\RateThirtyFive\ and \$\RateForty.}

So the Opus conclusion is a statement about Opus~5.5 \emph{at these prices}: if its cache reads cost about
\BEReadMultiple\X\ as much, routing sticky to Sonnet would roughly break even on this workload.

\begin{figure}[htbp]
\centering
% How a prompt cache behaves over a conversation: a schematic (no data).
\begin{tikzpicture}[x=1cm, y=1cm, font=\small,
  seg/.style={draw=black!60, minimum height=0.55cm, inner sep=0pt},
  wr/.style={seg, fill=okorange!75, postaction={pattern=north east lines, pattern color=white}},
  rd/.style={seg, fill=okblue!35},
  lab/.style={anchor=east, align=right, font=\small},
  note/.style={anchor=west, align=left, font=\footnotesize, text=black!75}]
  % row 1: first request
  \node[lab] at (0,0) {Turn 1, first request};
  \node[wr, minimum width=3.0cm, anchor=west] (a1) at (0.2,0) {tools + system};
  \node[wr, minimum width=1.2cm, anchor=west] at (a1.east) {turn 1};
  \node[note] at (4.6,0) {everything is written to the cache (cold start)};
  % row 2: same model, next turn
  \node[lab] at (0,-0.9) {Turn 2, same model};
  \node[rd, minimum width=3.0cm, anchor=west] (b1) at (0.2,-0.9) {tools + system};
  \node[rd, minimum width=1.2cm, anchor=west] (b2) at (b1.east) {turn 1};
  \node[wr, minimum width=1.2cm, anchor=west] at (b2.east) {turn 2};
  \node[note] at (5.8,-0.9) {old part read cheaply, only the new turn is written};
  % row 3: switch model
  \node[lab] at (0,-1.8) {Turn 3, other model};
  \node[wr, minimum width=3.0cm, anchor=west] (c1) at (0.2,-1.8) {tools + system};
  \node[wr, minimum width=1.2cm, anchor=west] (c2) at (c1.east) {turn 1};
  \node[wr, minimum width=1.2cm, anchor=west] (c3) at (c2.east) {turn 2};
  \node[wr, minimum width=1.2cm, anchor=west] at (c3.east) {turn 3};
  \node[note] at (7.0,-1.8) {each model has its own cache:\\ the whole history is written again};
  % row 4: long gap
  \node[lab] at (0,-2.9) {After a \LongGapMin-minute gap};
  \node[wr, minimum width=3.0cm, anchor=west] (d1) at (0.2,-2.9) {tools + system};
  \node[wr, minimum width=1.2cm, anchor=west] (d2) at (d1.east) {turn 1};
  \node[wr, minimum width=1.2cm, anchor=west] (d3) at (d2.east) {turn 2};
  \node[wr, minimum width=1.2cm, anchor=west] at (d3.east) {turn 3};
  \node[note] at (7.0,-2.9) {entries expire after \CacheTTLMin\ idle minutes:\\ written again};
  % key
  \node[rd, minimum width=0.6cm, anchor=west] (k1) at (0.2,-4.0) {};
  \node[anchor=west, font=\footnotesize] at (k1.east) {\ read from cache (Opus: \$\PriceOpusRead\ per million tokens)};
  \node[wr, minimum width=0.6cm, anchor=west] (k2) at (6.6,-4.0) {};
  \node[anchor=west, font=\footnotesize] at (k2.east) {\ written to cache (Opus: \$\PriceOpusWrite)};
\end{tikzpicture}

\caption{How the prompt cache behaves over a conversation (schematic, not to scale).}
\label{fig:cache}
\howtoread{Each bar is one request; it always re-sends everything before it. Blue parts are read back from
the cache at the cheap rate; hatched orange parts are written at the expensive rate. Within one model only
the new turn is written. A switch to another model, or a pause longer than the cache's lifetime, makes the
whole history expensive again.}
\end{figure}

\paragraph{What a probe of the provider showed.} Before designing the campaign we sent \ProbeRequests\
deliberately constructed requests straight to the provider (cost \$\ProbeUSD) to pin down the rules that
the design depends on. \Cref{tab:probe} shows the evidence; each row quotes the tokens written and read by
one request.

\begin{table}[htbp]
\centering\small
\caption{The cache-semantics probe. ``P'' is a prompt prefix of \ProbePTokens\ tokens; ``P+X'' a longer
prompt of about \ProbeFullTokens\ tokens; ``Y'' a further \ProbeYTokens-token extension.}
\label{tab:probe}
\begin{tabular*}{\textwidth}{@{\extracolsep{\fill}}>{\raggedright\arraybackslash}p{0.3\textwidth} >{\raggedright\arraybackslash}p{0.38\textwidth} r r@{}}
\toprule
Question & What was sent & tokens written & tokens read \\
\midrule
Longer entry does not serve its own prefix & Write prompt P+X, then ask for P alone & 10{,}948 & 0 \\
 & Ask for P alone again & 0 & 10{,}948 \\
A growing conversation reuses the earlier entry & Write P+X, then send P+X+Y & 2{,}272 & 15{,}455 \\
Two API keys do not share & Write with key 1, same request with key 2 & 15{,}455 & 0 \\
Simultaneous identical requests both pay & Two at the same instant (each) & 15{,}455 & 0 \\
Reads keep an entry alive & Read at 4 min, then again at 7 min & 0 & 15{,}455 \\
 & Never read, asked again at 7 min & 15{,}455 & 0 \\
Each model has its own cache & Written on Sonnet, then sent to Opus & 15{,}454 & 0 \\
Sonnet caches per effort level & Effort medium written, then no effort & 15{,}464 & 0 \\
 & Same pair on Opus & 0 & 15{,}464 \\
\bottomrule
\end{tabular*}

\end{table}

\begin{keyidea}
Six facts from the probe shaped the experiment:
\begin{enumerate}[leftmargin=1.4em]
  \item \textbf{A cache entry exists only where a request marked one, for exactly that prefix.} So one unique
    string at the very start of a session's system prompt (a \term{nonce}) makes every cache entry of that
    session private to it.
  \item \textbf{Each model has its own cache}, and on Sonnet each effort level has its own. A switch starts
    cold, whatever the timing: that cold start \emph{is} the switching cost we want to measure.
  \item \textbf{Two API keys do not share a cache}, but there were not enough keys for one per arm.
  \item \textbf{Two identical requests sent at the same instant both pay the write.} The nonce avoids this
    too, because no two sessions ever send identical requests.
  \item \textbf{Every read keeps an entry alive.} Staggering the arms in time would not isolate them, because
    a running session keeps its shared prefix warm.
  \item \textbf{Entries expire after \CacheTTLMin\ idle minutes.} Long pauses in a conversation are a
    real cost and must be part of the scenarios.
\end{enumerate}
\end{keyidea}
